\section{Balance slices with sharp dimension and barrier bounds}
\label{sec:balance-slices}

Connected extreme points force an extra cone coordinate under a strict
dimension cap, but they do not force one extra coordinate per factor.
We construct slices attaining the former bound for arbitrarily many
factors.  Their intrinsic barrier parameters are also exact.  This gives
an explicit separation between dimension resources and barrier resources.

Let $V_i$, $1\leq i\leq L$, be simple Euclidean Jordan algebras of
rank $r_i\geq2$, with cones of squares $K_i$, identities $e_i$, and
Jordan frames $(c_{ij})_{j=1}^{r_i}$.  Here $L\geq2$.
Choose a nonzero off-diagonal Peirce element
$A_i\in V_i(c_{i1},c_{i2})$, and put
\[
 \ell_i(X_i)=\ip{A_i}{X_i},\qquad
 M=\sum_i\dim V_i,\qquad \rho=\sum_i r_i.
\]
Define the balance cone and its trace-one base by
\begin{equation}\label{eq:balance-slice}
 \widehat Z=\left\{X\in\prod_iK_i:\sum_i\ell_i(X_i)=0\right\},
 \qquad
 Z=\left\{X\in\widehat Z:\sum_i\tr X_i=1\right\}.
\end{equation}
All interiors and barriers below are taken in the respective affine hulls.

\begin{theorem}[Balance slices]\label{thm:balance-slices}
The set $Z$ is a compact convex body of dimension $M-2$, with a
path-connected extreme-point set.  Each displayed factor is
indecomposable, its membership constraint is essential, and it occurs
in a nonconstant supporting slack.  For arbitrary self-concordant
barriers, including barriers coupling all blocks,
\begin{equation}\label{eq:balance-parameters}
 \vartheta_{\mathrm{opt}}(\ri\widehat Z)=\rho,
 \qquad
 \vartheta_{\mathrm{opt}}(\ri Z)=\rho-1.
\end{equation}
The product Jordan log-determinant, restricted to these domains, attains
both values.  Every closed convex affine lift of $Z$ with nonempty
bounded fibers has intrinsic barrier parameter at least $\rho-1$.
The conclusions include heterogeneous mixtures of classical Hermitian,
spin, and Albert factors.
\end{theorem}

The elementary minimal-face criterion used below is classical in finite
dimensions; see also the general injectivity formulation of
Henrion, Kru\v{z}\'ik, and Weis~\cite[Theorem~2.6]{HKW2026}.
The ambient symmetric-cone barrier theory is standard~\cite{FK1994,GT1998}.
The content of the theorem is the combination of the explicit balance
construction, connected extreme points, and exact parameters on its two
sections, rather than a new general extremality or barrier theorem.

\begin{proof}
The point $X_i^0=e_i/\rho$ is strictly feasible.  The total trace and
balance functionals are independent, so $\dim Z=M-2$.  Positivity
and fixed total trace imply compactness.

We first classify and connect the extreme points.  If a point $X$ belongs
to an affine section $\{X\in K:\mathcal A X=b\}$, it is extreme precisely
when
\begin{equation}\label{eq:balance-face-criterion}
 \ker\mathcal A\cap\spanop F_K(X)=\{0\}.
\end{equation}
Indeed, $X$ is relatively interior to its minimal face, so a nonzero
vector in this intersection gives a feasible segment through $X$.
Conversely, the endpoints of any such segment lie in the minimal face.
A rank-$q\geq2$ face in a simple Jordan cone has dimension
$q+a_iq(q-1)/2\geq3$, where $a_i$ is the Peirce multiplicity.
There are only two affine equations in~\eqref{eq:balance-slice}.
Thus each extreme point has one or two nonzero blocks, each on a
primitive ray.  Normalize its primitive idempotents to trace one, and
write $h_i(P)=\ell_i(P)$.  The constraint columns are $(1,h_i(P))$.
The exact possibilities are consequently
\begin{itemize}
 \item one block $P$ with $h_i(P)=0$;
 \item two blocks $wP$ and $(1-w)Q$ with $h_i(P)h_j(Q)<0$,
       where $w=-h_j(Q)/(h_i(P)-h_j(Q))$.
\end{itemize}
In the second case $i\ne j$.  These conditions also suffice by
\eqref{eq:balance-face-criterion}.

We record the sign geometry of $h_i$.  By a Jordan automorphism and
positive rescaling, its defining element becomes $c_{i1}-c_{i2}$;
this follows from the spectral theorem and transitivity on Jordan
frames.  For a classical factor $H_r(\mathbb F)$,
$\mathbb F\in\{\R,\C,\HH\}$, primitive idempotents are
$uu^*/\norm{u}^2$, and their height is proportional to
\[
 \frac{|u_1|^2-|u_2|^2}{\norm{u}^2}.
\]
The positive region, in the chart $u_1=1$, is
$\{u_2:|u_2|<1\}\times\mathbb F^{r-2}$ and is path-connected;
the negative region is obtained by interchanging coordinates.
If $r\geq3$, the zero region is path-connected: for $u_1\ne0$ it is
$\Sph^{\dim_\R\mathbb F-1}\times\mathbb F^{r-2}$, and every such
point can be joined, by letting a third coordinate tend to infinity,
to a fixed ray in the subspace $u_1=u_2=0$.
That subspace has connected projective space (or is one point).
This also joins the two chart components when $\mathbb F=\R$.
For $r=2$ the zero region is
$\Sph^{\dim_\R\mathbb F-1}$; only the real case has two points.
Every zero ray admits a path entering either strict sign region:
vary $|u_2|$ in the first chart, or introduce a first or second
coordinate at a ray with $u_1=u_2=0$.
For a spin factor, primitive idempotents form a sphere and height is
a nonzero linear coordinate.  Its sign regions are hemispheres, its
zero region is an equator, and the same assertions hold, with a
two-point equator in the three-dimensional cone.

For completeness the Albert case can be checked in coordinates,
without an analogy with associative projective spaces.
The chart of the octonionic projective plane with first diagonal
coordinate positive is parametrized by $(x,y)\in\OO^2$.
Its primitive idempotent is
\begin{equation}\label{eq:balance-albert-chart}
 P(x,y)=\frac{1}{1+|x|^2+|y|^2}
 \begin{pmatrix}
 1&\bar x&\bar y\\
 x&|x|^2&x\bar y\\
 y&y\bar x&|y|^2
 \end{pmatrix},
 \qquad
 h(P)=\frac{1-|x|^2}{1+|x|^2+|y|^2}.
\end{equation}
These are the usual reduced homogeneous coordinates
\cite[Section~3]{HSV2009}; the primitive-idempotent model is described
in~\cite{Baez2002}.  The two octonions $x,y$ generate an associative
subalgebra, which directly verifies $P^2=P$ and $\tr P=1$.
The positive region is $\B^8\times\R^8$, and permutation of the first
two frame idempotents identifies the negative region with it.
The zero region in this chart is $\Sph^7\times\R^8$.
Outside the chart, zero height forces both first diagonal coordinates
to vanish, leaving only $c_{i3}$: positivity makes their Peirce
off-diagonal components vanish as well.  A path $(x,ty_0)$ with
$|x|=|y_0|=1$ tends to $c_{i3}$ as $t\to\infty$.
After reparametrization this is a continuous path on a closed interval,
so the entire zero region is path-connected.  Varying $|x|$ gives both
signs near every finite zero point; keeping $|x|<1$ or $|x|>1$ and
letting $|y|\to\infty$ gives such paths at $c_{i3}$.

These facts give paths of extreme points of $Z$ as follows.
From any two-block extreme point, move its positive-height ray within
the positive region to a path ending at height zero, while holding
the negative-height ray fixed.  The displayed balance weights converge
to one and zero, so this joins the point to a one-block extreme point.
For prescribed zero-height rays $P$ and $Q$ in different blocks,
choose continuous paths $P_t$ and $Q_t$ such that $P_0=P$, $Q_1=Q$,
$h(P_t)>0$ for $t>0$, and $h(Q_t)<0$ for $t<1$.
The positive path may be held constant after entering its sign region,
and similarly for the negative path.  The uniquely balanced pair
extends at $t=0,1$ to $P,Q$, respectively.  Finally, any two zero
rays in one block can be connected through a zero ray in another
block.  This last step uses $L\geq2$ and also handles each two-point
zero region.  Thus all extreme points lie in one path component.

We next compute barrier parameters.  Let
$\Phi(X)=-\sum_i\log\det X_i$.  The frame-diagonal section of
$\widehat Z$ is $\R_+^\rho$, because each $A_i$ is off diagonal.
It meets $\ri\widehat Z$, and its boundary belongs to
$\partial\widehat Z$.  To see the latter assertion, a vanished frame
coordinate is detected by a nonnegative primitive functional, which
is strictly positive at $X^0$; it is therefore a nonconstant supporting
functional on the balance section.  Its trace-one section is a simplex
of dimension $\rho-1$ with the same boundary inheritance.
Restriction of any barrier to these sections preserves its parameter.
The classical orthant and simplex lower bounds
\cite[Section~2.3.4 and Proposition~2.3.6]{NN1994} give respectively
$\rho$ and $\rho-1$, even for coupled barriers.

The standard barrier gives the upper bound $\rho$ on $\widehat Z$.
For the sharper normalized bound, let $g=\nabla\Phi$, $H=\nabla^2\Phi$,
and let $a$ be the total-trace covector.  Jordan identities give
\[
 H^{-1}g=-X,\qquad g^TH^{-1}g=\rho,\qquad
 a^TH^{-1}a=\sum_i\tr(X_i^2)\leq\left(\sum_i\tr X_i\right)^2.
\]
Here $H_i=P(X_i^{-1})$ and $H_i^{-1}=P(X_i)$, including in the
Albert algebra.  On the trace-one hyperplane the squared dual norm
of the restricted gradient is exactly
\begin{equation}\label{eq:balance-gradient}
 \rho-\frac{(a^TH^{-1}g)^2}{a^TH^{-1}a}
 =\rho-\frac{1}{\sum_i\tr(X_i^2)}\leq\rho-1.
\end{equation}
Restriction to the balance tangent can only decrease this norm.
Self-concordance and boundary divergence survive affine restriction,
which proves~\eqref{eq:balance-parameters}.
Lemma~\ref{lem:bounded-fiber-barrier} proves the assertion for other
closed lifts with bounded fibers; no restriction on cross-block
derivatives of their barriers is imposed.

Finally, simplicity makes each factor indecomposable.  If membership
in $K_i$ is removed, choose $j\ne i$ and set
$X_i=-Te_i/r_i$, $X_j=(T+1)e_j/r_j$, with all other blocks zero.
This is an unbounded family satisfying both affine equations and
all remaining cone constraints.  Each factor also supports a
nonconstant nonnegative slack: $X\mapsto\ip{c_{i2}}{X_i}$ vanishes
at the feasible one-block point $c_{i1}$ and is positive at $X^0$.
\end{proof}

\begin{corollary}[An exact strict-cap frontier]\label{cor:balance-cap}
Suppose the factors in~\eqref{eq:balance-slice} all have dimension $d$,
and put $N=dL-2$.  Among full-slack factorizations of $Z$ through products
of arbitrary nonzero proper cones of dimension at most $d$, the minimum
factor count and minimum total cone dimension are
\[
 L_{\min}=L,\qquad M_{\min}=N+2=dL.
\]
These statements permit row splitting and impose no regularity on the
factor maps.
\end{corollary}
\begin{proof}
For a competing product, ordinary slack rank gives $M'\geq N+1$.
Equality would identify that product with the homogenization cone of
$Z$ by Theorem~\ref{thm:minimal-dimension-rigidity}.
The projectivized extreme rays of a nontrivial product form a
disconnected union, whereas those of the homogenization are
homeomorphic to the connected set of extreme points of $Z$.
Equality therefore requires a single factor.  This violates the cap,
since $d<N+1=dL-1$.  Consequently $M'\geq N+2=dL$ and
$L'\geq M'/d\geq L$.  The displayed relative-Slater slice gives
an attaining full-slack factorization by Lemma~\ref{lem:reduction}.
\end{proof}

\begin{corollary}[An alternate classical moment]\label{cor:balance-alternate}
For $L$ copies of $H_r(\mathbb F)_+$, $r\geq2$, the preceding
theorem and corollary remain valid with
$A_i=\mathbf1\mathbf1^*-I_r$ instead of an off-diagonal Peirce element.
Writing $\delta=\dim_\R\mathbb F$, the zero-height primitive rays in
one block form $\Sph^{\delta(r-1)-1}$.
\end{corollary}
\begin{proof}
This matrix has zero diagonal and signature $(1,r-1)$, with eigenvalues
$r-1,-1,\ldots,-1$.  The Slater, barrier, and essentiality proofs
therefore apply unchanged.  In eigenvector coordinates $(a,b)$,
the moment is $((r-1)|a|^2-\norm b^2)/(|a|^2+\norm b^2)$.
A zero ray has the unique representative
$(1,\sqrt{r-1}\,u)$ with $\norm u=1$, proving the sphere assertion.
Every strict-sign ray can be moved to a zero ray while retaining
its sign until the endpoint, by scaling the two components; if one
component vanishes, first introduce it.  Two prescribed zero rays in
distinct blocks admit an explicit extreme path.  For their fixed
unit vectors $u$, replace the representative by
$(1,\lambda(c)\sqrt{r-1}\,u)$, where
\[
 \lambda(c)^2=\frac{r-1-c}{(r-1)(1+c)},\qquad -1<c<r-1.
\]
Fix $0<c_0<1$ and assign moments $tc_0$ and $-(1-t)c_0$.
Their balance weights are $1-t$ and $t$.  This connects the two
prescribed one-block points.  A second block again connects the
two components when $\mathbb F=\R$, $r=2$.
Connectedness and the cap conclusion follow exactly as above.
\end{proof}

At matched cone dimension $d=r+\delta r(r-1)/2$, the classical family
and the spin-factor family built from $Q_d^L$ both have body dimension
$dL-2$, minimum capped factor count $L$, and minimum total dimension
$dL$.  Their optimal direct-body barrier parameters are $rL-1$ and
$2L-1$.  These are two different bodies.  The comparison proves that
count and dimension do not determine intrinsic barrier cost; it does
not construct a Lorentz lift of the classical body.
For Albert factors the same formulas give $N=27L-2$, $M_{\min}=27L$,
and intrinsic parameter $3L-1$.

The sharp family-specific combination above follows from classical
minimal-face, Jordan-barrier, and simplex arguments, together with the
explicit connectivity proof and minimum-dimension rigidity.
No priority claim is attached to the alternate moment construction.
The relation to intrinsic path length requires an additional metric
calculation and is treated separately below.

\begin{example}[A failed coupled logarithmic cancellation]
In two three-dimensional Lorentz factors write
$q_i=t_i^2-x_i^2-y_i^2$, and impose $x_1+x_2=0$ and
$t_1+t_2=1$. The candidate
$\Psi=-\log q_1-\log q_2+\log(q_1+q_2)$ cancels a boundary
logarithmic order, but is not self-concordant. At
$(t_1,t_2,x_1,x_2,y_1,y_2)=(1/3,2/3,0,0,0,1/3)$, the feasible
direction varying only $y_2$ gives $D^2\Psi=13/4$ and
$D^3\Psi=25$. Thus $(D^3\Psi)^2>4(D^2\Psi)^3$.
For example these derivatives follow directly from
$q_1=1/9$, $q_2=1/3-2h/3-h^2$ along that line.
This illustrates why the parameter proof above uses the restricted
Hessian metric and an independent lower section, not boundary orders alone.
\end{example}
